% ==============================================================================
% TCC - Nome do Aluno
% Capítulo 1 - Introdução
% ==============================================================================
\chapter{Introdução}
\label{sec-intro}

Sistemas de \textit{Retrieval-Augmented Generation} (RAG) tornaram-se uma das
abordagens mais utilizadas para permitir que modelos de linguagem de grande
porte (\textit{Large Language Models} -- LLMs) respondam perguntas com base em
uma base de conhecimento externa e específica de domínio, reduzindo problemas
como respostas desatualizadas ou alucinadas \cite{lewis2020retrieval}. Na abordagem tradicional de RAG, o texto é dividido em fragmentos
(\textit{chunks}), convertido em vetores por meio de \textit{embeddings} e
recuperado por similaridade semântica no momento da consulta. Essa estratégia,
no entanto, apresenta limitações quando o domínio de conhecimento possui uma
estrutura fortemente relacional, na qual a resposta correta depende de conectar
informações dispersas em diferentes trechos do texto -- situação recorrente em
textos legislativos, nos quais artigos remetem uns aos outros, estabelecem
exceções e são complementados por súmulas e jurisprudência.

O \textit{framework} GraphRAG, desenvolvido pela Microsoft \cite{edge2024global}, propõe uma alternativa a essa limitação: em vez de recuperar apenas
fragmentos de texto semanticamente similares, o GraphRAG constrói um grafo de
conhecimento a partir do corpus de entrada, extraindo entidades e relações por
meio de um LLM, agrupando-as em comunidades hierárquicas e gerando resumos para
essas comunidades. Esse processo permite tanto buscas locais, focadas em
entidades específicas, quanto buscas globais, capazes de sintetizar informação
distribuída por múltiplas partes do corpus -- uma capacidade particularmente
relevante para a análise de legislação.

Este trabalho tem como contexto a aplicação do GraphRAG à legislação do Código
de Defesa do Consumidor (CDC, Lei n.º 8.078/1990). O escopo do trabalho está
delimitado à construção de uma prova de conceito (\textit{Proof of Concept} --
PoC): um sistema que indexa o texto do CDC utilizando o GraphRAG, armazenando
as informações como um grafo de entidades e relações juridicamente
relevantes, e que disponibiliza uma aplicação web para (i) visualização
interativa do grafo de conhecimento construído e (ii) realização de consultas
em linguagem natural, utilizando as diferentes modalidades de busca oferecidas
pelo GraphRAG (local, global, \textit{drift} e \textit{basic}). Ficam fora do
escopo deste trabalho a extensão do sistema a outras legislações e a disponibilização do
sistema em ambiente de produção -- delimitações necessárias para viabilizar a
conclusão do trabalho dentro do tempo hábil de um Trabalho de Conclusão de
Curso.

%%% Início de seção. %%%
\section{Motivação e Justificativa}
\label{sec-intro-motjus}

A motivação deste trabalho parte do interesse em investigar, na prática, as
limitações dos sistemas de RAG tradicionais ao lidar com textos que possuem
forte densidade de relações internas, como é o caso da legislação. Textos
jurídicos raramente podem ser interpretados de forma isolada: um único artigo
frequentemente depende de definições estabelecidas em outros artigos, está
sujeito a exceções previstas em dispositivos distintos e é interpretado à luz
de súmulas e entendimentos consolidados. Uma abordagem de RAG baseada apenas em
similaridade vetorial entre fragmentos de texto tende a perder essas conexões,
recuperando trechos isolados que, embora semanticamente relacionados à
pergunta, não contemplam a cadeia de relações necessária para uma resposta
completa e juridicamente coerente.

O Código de Defesa do Consumidor foi escolhido como domínio de aplicação por
três razões principais. Primeiro, trata-se de uma legislação de grande relevância 
social e de ampla aplicação no cotidiano da sociedade brasileira, tornando-se um domínio de interesse para a investigação e avaliação da abordagem proposta. Segundo, o CDC possui uma estrutura de
tamanho gerenciável para os fins de uma PoC, permitindo iterações de
desenvolvimento e ajuste do sistema em tempo hábil, ao mesmo tempo em que
preserva a complexidade relacional típica de textos legislativos. Terceiro, o crescimento do
interesse em soluções de acesso à justiça apoiadas por inteligência
artificial reforça a relevância prática de investigar sistemas de recuperação
de informação mais adequados a esse domínio.

Ainda que sistemas de RAG aplicados ao domínio jurídico já tenham sido
explorados na literatura, a maior parte desses trabalhos utiliza abordagens de
recuperação vetorial tradicional. Este trabalho busca
contribuir com essa lacuna, investigando de forma prática se a abordagem em
grafo do GraphRAG traz ganhos perceptíveis de qualidade e completude nas
respostas, quando comparada à abordagem vetorial tradicional, em um domínio
legislativo de escopo delimitado.

%%% Início de seção. %%%
\section{Objetivos}
\label{sec-intro-obj}

\subsection*{Objetivo Geral}

Desenvolver uma prova de conceito de um sistema de \textit{Retrieval-Augmented
Generation} baseado em grafo de conhecimento (GraphRAG) para indexação e
consulta da legislação do Código de Defesa do Consumidor, disponibilizando uma
aplicação web que permita a visualização do grafo de conhecimento construído e
a realização de consultas em diferentes modalidades de busca.

\subsection*{Objetivos Específicos}

\begin{itemize}
    \item Revisar a literatura sobre sistemas de \textit{Retrieval-Augmented
    Generation}, RAG baseado em grafos de conhecimento e aplicações de
    Processamento de Linguagem Natural no domínio jurídico.

    \item Definir uma taxonomia de entidades e relações adequada à estrutura
    do Código de Defesa do Consumidor, a ser utilizada no processo de extração
    de grafo.

    \item Implementar o \textit{pipeline} de indexação do GraphRAG sobre o
    texto do CDC, compreendendo a extração de entidades e relações, a
    construção do grafo e a detecção de comunidades.

    \item Desenvolver uma aplicação web para visualização interativa do grafo
    de conhecimento gerado e realização de consultas em linguagem natural, permitindo
    a utilização das diferentes modalidades de busca disponibilizadas pelo GraphRAG 
    (\textit{local}, \textit{global}, \textit{drift} e \textit{basic}).

    \item Avaliar comparativamente a qualidade e a completude das respostas
    obtidas por meio do GraphRAG em relação a uma abordagem de RAG vetorial
    tradicional, utilizada como \textit{baseline}.
\end{itemize}

%%% Início de seção. %%%
\section{Método de Desenvolvimento do Trabalho}
\label{sec-intro-met}

O desenvolvimento deste trabalho caracteriza-se como uma pesquisa aplicada, de
natureza tecnológica, cujo produto final é uma prova de conceito funcional. O
trabalho será conduzido por meio das seguintes etapas metodológicas:

\begin{itemize}
    \item \textbf{Revisão bibliográfica:} levantamento de referenciais
    teóricos sobre RAG, RAG baseado em grafos (com foco no \textit{framework}
    GraphRAG) e extração de entidades e relações via LLMs.

    \item \textbf{Preparação do corpus:} obtenção e pré-processamento do texto
    do Código de Defesa do Consumidor, incluindo a definição da granularidade
    de segmentação (\textit{chunking}) do texto para a etapa de indexação.

    \item \textbf{Definição da taxonomia do domínio:} especificação dos tipos
    de entidade e relação relevantes para a legislação em análise (por
    exemplo, artigo, direito, obrigação, sanção, relação de consumo), a serem
    utilizados para orientar a extração de grafo realizada pelo LLM.

    \item \textbf{Configuração do ambiente de indexação:} montagem do
    ambiente de execução do GraphRAG, incluindo a definição do(s) modelo(s) de
    linguagem utilizados para as etapas de extração, resumo e geração de
    respostas, executados em ambiente conteinerizado para garantir
    reprodutibilidade.

    \item \textbf{Execução e ajuste do \textit{pipeline} de indexação:}
    execução do \textit{pipeline} de indexação do GraphRAG sobre o corpus,
    seguida de validação manual da qualidade das entidades e relações
    extraídas, com ajustes iterativos de configuração e de \textit{prompts}
    até que a extração atinja um nível satisfatório de cobertura e
    coerência.

    \item \textbf{Desenvolvimento da aplicação web:} implementação de um
    \textit{backend} responsável por expor as funcionalidades de consulta do
    GraphRAG (busca local, global, \textit{drift} e \textit{basic}) e de um
    \textit{frontend} para visualização interativa do grafo de conhecimento e
    para a submissão de consultas em linguagem natural.

    \item \textbf{Avaliação:} definição de um conjunto de perguntas de teste
    sobre o CDC e comparação qualitativa das respostas obtidas pelas
    diferentes modalidades de busca do GraphRAG frente a uma abordagem de RAG
    vetorial tradicional, utilizada como \textit{baseline} de comparação.
\end{itemize}

%%% Início de seção. %%%
\section{Cronograma}
\label{sec-intro-crono}

As atividades previstas para o desenvolvimento deste trabalho, alinhadas aos
objetivos específicos apresentados na Seção~\ref{sec-intro-obj}, são descritas
a seguir:

\begin{itemize}
    \item \textbf{Atividade 1 -- Revisão bibliográfica:} levantamento e
    leitura de referenciais teóricos sobre RAG, GraphRAG e PLN aplicada ao
    domínio jurídico, culminando na escrita do referencial teórico do
    trabalho.

    \item \textbf{Atividade 2 -- Definição da taxonomia e preparação do
    corpus:} definição dos tipos de entidade e relação do domínio, e
    preparação do texto do CDC para indexação.

    \item \textbf{Atividade 3 -- Configuração do ambiente de indexação:}
    montagem do ambiente conteinerizado e definição dos modelos de linguagem
    utilizados.

    \item \textbf{Atividade 4 -- Execução e ajuste da indexação:} execução
    iterativa do \textit{pipeline} de indexação do GraphRAG, com validação e
    ajuste da qualidade da extração do grafo.

    \item \textbf{Atividade 5 -- Desenvolvimento do \textit{backend}:}
    implementação da camada de \textit{backend} responsável por expor as
    consultas ao GraphRAG.

    \item \textbf{Atividade 6 -- Desenvolvimento do \textit{frontend}:}
    implementação da interface web de visualização do grafo e de consulta.

    \item \textbf{Atividade 7 -- Avaliação comparativa:} definição do
    conjunto de perguntas de teste e comparação entre as modalidades de busca
    do GraphRAG e o \textit{baseline} vetorial.

    \item \textbf{Atividade 8 -- Escrita e revisão do TCC:} consolidação dos
    resultados e redação final do documento do trabalho.
\end{itemize}

\begin{table}[h]
\centering
\caption{Cronograma de execução das atividades}
\label{tab-cronograma}
\begin{tabular}{c|c|c|c|c|c|c|}
\cline{2-7}
\multicolumn{1}{l|}{} & \textbf{Mês 1} & \textbf{Mês 2} & \textbf{Mês 3} & \textbf{Mês 4} & \textbf{Mês 5} & \textbf{Mês 6} \\ \hline
\multicolumn{1}{|c|}{\textbf{Atividade 1}} & X & X &  &  &  &  \\ \hline
\multicolumn{1}{|c|}{\textbf{Atividade 2}} &  & X & X &  &  &  \\ \hline
\multicolumn{1}{|c|}{\textbf{Atividade 3}} &  & X &  &  &  &  \\ \hline
\multicolumn{1}{|c|}{\textbf{Atividade 4}} &  & X & X & X &  &  \\ \hline
\multicolumn{1}{|c|}{\textbf{Atividade 5}} &  &  &  & X & X &  \\ \hline
\multicolumn{1}{|c|}{\textbf{Atividade 6}} &  &  &  & X & &  \\ \hline
\multicolumn{1}{|c|}{\textbf{Atividade 7}} &  &  &  &  & X & X \\ \hline
\multicolumn{1}{|c|}{\textbf{Atividade 8}} & X & X & X & X & X & X \\ \hline
\end{tabular}
\end{table}